\chapter{微分流形}
    % \section{Kunn\"{e}th 公式与 Leray-Hirsch 定理}
    % \begin{theorem}[Kunn\"{e}th 公式]
    %     设流形 $M$ 有有限好覆盖, $F$ 是任意流形, 则
    %     \begin{equation*}
    %         H^*(M\times F)\cong H^*(M)\otimes H^*(F)
    %     \end{equation*}
    % \end{theorem}

    % \begin{proof}[证明概要]
    %     设 $\pi:M\times F\rightarrow M$, $\rho:M\times F\rightarrow F$ 
    %     为乘积流形到两个分量的投影, 则可以定义
    %     \begin{align*}
    %         \psi:H^*(M)\otimes H^*(F)&\rightarrow H^*(M\times F) \\
    %         \omega\otimes\tau&\mapsto\pi^*\omega\wedge\rho^*\tau
    %     \end{align*}
    %     由 M-V 论证, 可以得到如下交换图:

    %     \begin{tikzcd}[column sep = tiny]
    %         \vdots \arrow[d]                                                                                                 & \vdots \arrow[d]                                          \\
    %         \bigoplus\limits_{p+q=k}H^p(U\cup V)\otimes H^q(F) \arrow[d] \arrow[r, "\psi"]                                   & H^k\left((U\cup V)\times F\right) \arrow[d]               \\
    %         \bigoplus\limits_{p+q=k}\left((H^p(U)\otimes H^q(F))\oplus (H^p(V)\otimes H^q(F))\right) \arrow[r, "\psi"] \arrow[d] &             H^k(U\times F)\oplus H^k(V\times F) \arrow[d] \\
    %         \bigoplus\limits_{p+q=k}H^p(U\cap V)\otimes H^q(F) \arrow[r, "\psi"] \arrow[d, "\md^*"]                          & H^k((U\cap V)\times F) \arrow[d, "\md^*"]                 \\
    %         \bigoplus\limits_{p+q=k+1}H^p(U\cup V)\otimes H^q(F) \arrow[d] \arrow[r, "\psi"]                                 & H^{k+1}\left((U\cup V)\times F\right) \arrow[d]           \\
    %         \vdots                                                                                                           & \vdots                                                   
    %     \end{tikzcd}

    %     前两个圈的交换性显然, 第三个圈的交换性需要用到 $\md^*$ 的表达式. 
    %     由五引理能得到归纳递推, 归纳奠基是平凡的.
    % \end{proof}

    % \begin{theorem}[Leray-Hirsch 定理]
    %     设 $\pi:E\rightarrow M$ 为纤维丛, 纤维为 $F$, 
    %     若存在 $E$ 上的微分形式 $\{\ft{e}{1}{n}\}$ 满足将它们限制在每个纤维 
    %     $F_x$ 上都能得到 $H^*(F_x)$ 的一组基, 
    %     则
    %     \begin{equation*}
    %         H^*(E)\cong H^*(M)\otimes\{\ft{e}{1}{n}\}
    %         \cong H^*(M)\otimes H^*(F).
    %     \end{equation*}
    % \end{theorem}
    % \begin{proof}[证明概要]
    %     这里的关键在于不存在 $E$ 到 $F$ 的整体投影 $\rho$, 
    %     也就无法通过这个方式定义 $\rho^*:H^*(F)\rightarrow H^*(E)$ 了. 
    %     但是借助 $\{\ft{e}{1}{n}\}$ 我们可以构造合适的映射 $\tilde{\rho^*}$, 
    %     做法如下:固定某个点 $x\in M$, 也即固定某个纤维 $F_x$, 
    %     取 $H^*(F_x)$ 的一组基 $\{\ft{f}{1}{n}\}$, 定义:
    %     \begin{align*}
    %         \tilde{\rho^*}:H^*(F)&\rightarrow H^*(E) \\
    %         \sum_{i}a_if_i&\mapsto\sum_{i}a_ie_i.
    %     \end{align*} 

    %     于是可以定义:
    %     \begin{align*}
    %         \tilde{\psi}:H^*(M)\otimes H^*(F)&\rightarrow H^*(E) \\
    %         \omega\otimes\tau&\mapsto\pi^*\omega\wedge\tilde{\rho^*}\tau
    %     \end{align*}
    %     归纳递推仍由 M-V 论证给出; 
        
    %     \begin{tikzcd}[column sep = tiny]
    %         \vdots \arrow[d]                                                                                                 & \vdots \arrow[d]                                          \\
    %         \bigoplus\limits_{p+q=k}H^p(U\cup V)\otimes H^q(F) \arrow[d] \arrow[r, "\psi"]                                   & H^k\left(\pi^{-1}(U\cup V)\right) \arrow[d]               \\
    %         \bigoplus\limits_{p+q=k}\left((H^p(U)\otimes H^q(F))\oplus (H^p(V)\otimes H^q(F))\right) \arrow[r, "\psi"] \arrow[d] &             H^k(\pi^{-1}(U))\oplus H^k(\pi^{-1}(V)) \arrow[d] \\
    %         \bigoplus\limits_{p+q=k}H^p(U\cap V)\otimes H^q(F) \arrow[r, "\psi"] \arrow[d, "\md^*"]                          & H^k(\pi^{-1}(U\cap V)) \arrow[d, "\md^*"]                 \\
    %         \bigoplus\limits_{p+q=k+1}H^p(U\cup V)\otimes H^q(F) \arrow[d] \arrow[r, "\psi"]                                 & H^{k+1}\left(\pi^{-1}(U\cup V)\right) \arrow[d]           \\
    %         \vdots                                                                                                           & \vdots                                                   
    %     \end{tikzcd}
        
    %     因为 good cover 中的每个开集都同伦于单点, 
    %     而同伦映射诱导同构的拉回丛(注意这里是纤维丛的版本), 
    %     因此 good cover 同时也是 locally trivialization.
    %     故 $\pi^{-1}(U_\alpha)\cong U_{\alpha}\times F$, 
    %     从而 $H^*(\pi^{-1}(U_\alpha))\cong H^*(U_{\alpha})\otimes H^*(F)$, 
    %     这给出了归纳奠基.
    % \end{proof}
    % \begin{remark}
    %     实际上当底空间 $M$ 连通时, 定理的条件可弱化为 $\{\ft{e}{1}{n}\}$ 
    %     限制在某个纤维 $F_x$ 上得到 $H^*(F_x)$ 的一组基. 
    %     因为对不同的两点 $x,y$, 有道路 $\gamma:[0,1]\rightarrow M$ 将他们相连, 
    %     于是嵌入映射 
    %     $\iota_x:F_x\hookrightarrow E$ 和 $\iota_y:F_y\hookrightarrow E$ 
    %     同伦.
    %     因此拉回映射 $\iota^*_x:H^*(E)\rightarrow H^*(F_x)$ 
    %     与 $\iota^*_y:H^*(E)\rightarrow H^*(F_y)$ 相等.
    % \end{remark}
    % \begin{remark}
    %     这里的同构 $H^*(E)\cong H^*(M)\otimes H^*(F)$ 并不保持环结构(例如？)
    %     因此只能说 $H^*(E)$ 可看成 $H^*(M)$-模.
    % \end{remark}
    % \begin{remark}
    %     存在不满足{\rm Leray-Hirsch}定理条件的纤维丛, 比如{\rm Hopf}纤维化:
    %     \begin{center}
    %         \begin{tikzcd}
    %             S^1 \arrow[r] & S^3 \arrow[d] \\
    %                           & S^2          
    %         \end{tikzcd}    
    %     \end{center}
    %     其中
    %     \begin{equation*}
    %         H^*(S^3)\neq H^*(S^1)\otimes H^*(S^2)
    %     \end{equation*}
    % \end{remark}

    % \section{微分流形中的同伦}
    % \subsection{连续映射同伦于光滑映射}
    %     \begin{theorem}[用光滑映射逼近连续映射]
    %         设 $f:M^m\rightarrow\mR^n$ 为任意连续映射, 
    %         $\varepsilon:M\rightarrow\mR_{>0}$ 为任意误差函数.
    %         则存在光滑映射 $g:M\rightarrow\mR^n$ 满足 
    %         $||f(x)-g(x)||<\varepsilon(x),\;\forall\,x\in M$.
    %     \end{theorem}
    %     \begin{proof}
    %         对 $M$ 中任一点 $p$, 存在开邻域 $U_p$ 使得 
    %         \begin{equation*}
    %             ||f(q)-f(p)||<\varepsilon(p),\quad\varepsilon(q)
    %             >\frac{1}{2}\varepsilon(p),\quad\forall\,q\in U_p.
    %         \end{equation*}
    %         记 $\{\rho_p\}$ 为从属于开覆盖 $\{U_p\}$ 的单位分解, 定义 
    %         \begin{equation*}
    %             g(x) = \sum_{p}\rho_p(x)f(p) = \sum_{U_p\ni x}\rho_p(x)f(p).
    %         \end{equation*}
    %         则 $g(x)$ 是 $M$ 到 $\mR^n$ 的光滑映射且满足 
    %         \begin{align*}
    %             ||g(x)-f(x)|| &= \Big|\Big|\sum_{p}\rho_p(x)f(p) 
    %             - \sum_{p}\rho_p(x)f(x)\Big|\Big| \\
    %             &= \Big|\Big|\sum_{U_p\ni x}\rho_p(x)
    %             \Big(f(p)-f(x)\Big)\Big|\Big| \\
    %             &\leqslant\sum_{U_p\ni x}\rho_p(x)
    %             \big|\big|f(p)-f(x)\big|\big| \\
    %             &\leqslant\sum_{U_p\ni x}\rho_p(x)\varepsilon(x) \\
    %             &=\varepsilon(x).
    %         \end{align*}
    %         $g(x)$ 即为所求.
    %     \end{proof}
    %     \begin{remark}
    %         这个定理是单位分解的一个重要应用.
    %     \end{remark}
    %     \begin{remark}
    %         若 $f$ 的像集 $f(M)$ 包含于 $\mR^n$ 的某个开集 $U$ 中, 
    %         则也可使 $g(M)\subset U$. 具体做法是稍微修改一下误差函数.
    %         令 
    %         \begin{equation*}
    %             \varepsilon_1(x):=\min\left\{
    %                 \varepsilon(x),\frac{1}{2}d(f(x),\mR^n\backslash U)
    %                 \right\}
    %         \end{equation*} 
    %         以新误差函数 $\varepsilon_1$ 构造的 $g$ 自动满足 $g(M)\in U$.
    %     \end{remark}
    %     \begin{theorem}[用光滑映射逼近连续映射且保持在某个闭集不动]
    %         设 $f:M^m\rightarrow\mR^n$ 为连续映射, 
    %         且在闭子集 $A$ 的某个开邻域 $U_0$ 内光滑, 
    %         则对任意误差函数 $\varepsilon:M\rightarrow\mR_{>0}$,
    %         存在光滑映射 $g:M\rightarrow\mR^n$ 满足 
    %         $||f(x)-g(x)||<\varepsilon(x),\;\forall\,x\in M$ 且 
    %         $f\big|_A = g\big|_A$.
    %     \end{theorem}
    %     \begin{proof}
    %         主要的构造思路保持不变, 但开覆盖有些许变化. 
    %         对任意 $p\in M\backslash A$, 
    %         取开邻域 $U_p\subset M\backslash A$ 使得 
    %         \begin{equation*}
    %             ||f(q)-f(p)||<\varepsilon(p),\quad\varepsilon(q)
    %             >\frac{1}{2}\varepsilon(p),\quad\forall\,q\in U_p.
    %         \end{equation*}
    %         设 $\{\rho_0\}\cup\{\rho_p\}_{p\notin A}$ 
    %         为从属于开覆盖 $\{U_0\}\cup\{U_p\}_{p\notin A}$ 的单位分解, 令
    %         \begin{equation*}
    %             g(x) = \rho_0(x)f(x) + \sum_{p\notin A}\rho_p(x)f(p).
    %         \end{equation*}
    %         则因为
    %         \begin{equation*}
    %             f(x) = \rho_0(x)f(x) + \sum_{p\notin A}\rho_p(x)f(x), 
    %         \end{equation*}
    %         可以验证 $g$ 满足定理要求.
    %     \end{proof}
    %     \begin{remark}
    %         同上一个注一样, 该定理可以加强使逼近前后的像集在同一个开集中.
    %     \end{remark}
    %     \begin{theorem}[任意连续映射都同伦于一个光滑映射]
    %         任意光滑流形间的连续映射 $f:M^m\rightarrow N^n$ 
    %         均同伦于某一个光滑映射 $g:M\rightarrow N$.
    %     \end{theorem}
    %     \begin{proof}
    %         $f$ 在局部上可以看成是 $M$ 到 $\mR^n$ 的连续映射, 
    %         用之前的定理可以在局部上用光滑映射逼近 $f$, 
    %         则局部上可以用直线同伦连接 $f$ 与光滑映射, 
    %         对每个点的局部都如此考虑即可.

            
    %     \end{proof}

    % \section{Thom同构、Thom空间与Thom类}
    %     设 $\pi:E\rightarrow M$ 是 $M$ 上的秩为 $k$ 的可定向向量丛, 
    %     在 {\rm Bott \& Tu } 的书中我们知道存在 $\Phi\in H^k_{cv}(E)$ 使得 
    %     \begin{align*}
    %         \mathcal{T}:H^{*-k}(M)&\rightarrow H^*_{cv}(E) \\
    %         \omega&\mapsto\pi^*\omega\wedge\Phi
    %     \end{align*}
    %     是同构. 下面我们将用相对上同调以及 Thom 空间的语言重新描述 Thom 同构. 

    %     由定理\,\ref{thm:c-cohomology_v.s._reductive_cohomology}, 
    % 
    \section{主丛、主丛上的联络}
        \begin{definition}[$G$-主丛]
            设 $M$ 为流形, $G$ 为李群. 称流形 $P$ 为 $M$ 上的 $G$-主丛, 若
            \begin{itemize}
                \item $P$ 上有 $G$ 的自由右作用 
                $\mu:P\times G\rightarrow P$, $\mu(p,g)=p\cdot g$;
                \item $M$ 微分同胚于商空间 $P/G$, 
                且投影 $\pi:P\rightarrow M$ 是光滑的;
                \item $P$ 在 $M$ 上局部平凡, 即对任意 $x\in M$, 
                存在 $x$ 的开邻域 $U$ 及 $G$-等变的微分同胚 
                \begin{align*}
                    \psi:\pi^{-1}(U) & \rightarrow U\times G \\
                    p &\mapsto (\pi(p),\varphi(p)).
                \end{align*}
                其中 $G$ 在 $U\times G$ 上的右作用为 
                \begin{equation*}
                    (x,h)\cdot g = (x,hg), \quad\forall\,x\in U,\,h,g\in G,
                \end{equation*}
                因此 $G$-等变性即 
                \begin{equation*}
                    \psi(p\cdot g) = (\pi(p),\varphi(p\cdot g))
                    = (\pi(p),\varphi(p)g) = \psi(p)\cdot g. 
                \end{equation*}
            \end{itemize}
        \end{definition}
        在上述定义中我们强调了李群 $G$ 的右作用, 当然也有其它定义方式: 
        \begin{definition}[纤维丛, $G$-主丛]
            设 $M,E,F$ 均为流形, 李群 $G$ 左作用于 $F$. 
            称淹没 $\pi:E\rightarrow M$ 为以 $F$ 为纤维,
            $G$ 为结构群的纤维丛, 若对任意 $x\in M$,
            存在 $x$ 的开邻域 $U$ 及微分同胚(称为局部平凡化)
            \begin{align*}
                \psi:\pi^{-1}(U) & \rightarrow U\times F \\
                p &\mapsto (\pi(p),\varphi(p)).
            \end{align*}
            且当 $U_\alpha\cap U_\beta\neq\varnothing$ 时, 
            \begin{align*}
                \psi_\beta\circ\psi_\alpha^{-1}:(U_\alpha\cap U_\beta)
                \times F & \rightarrow (U_\alpha\cap U_\beta)\times F \\
                (x,f) & \mapsto (x,g_{\alpha\beta}(x)\cdot f).
            \end{align*}
            其中 $g_{\alpha\beta}:U_\alpha\cap U_\beta\rightarrow G$ 
            为光滑映射(称其为转移函数).

            特别地, 当纤维 $F$ 就是群 $G$ 且 $G$ 左作用于自身为左乘时, 
            称 $E$ 为 $M$ 上的 $G$-主丛.
        \end{definition}
        \begin{remark}
            两个定义是等价的, 从第一个定义推第二个定义的关键是得到转移函数, 
            而这是由于 $G$ 到自身的右 $G$-等变微分同胚只有左乘; 
            从第二个定义推第一个定义的关键是构造 $G$ 在全空间上的右作用, 
            这可以通过局部平凡化构造, 而转移函数保证了构造的良定义性.
        \end{remark}

        \begin{example}
            \begin{itemize}
                \item $M\times G$ 是 $M$ 上的平凡 $G$-主丛;
                \item $S^1$ 通过复数的乘法作用在 $S^{2n+1}$ 上, 
                商空间是 $\mC P^n$. 于是 $\pi:S^{2n+1}\rightarrow\mc P^n$ 
                给出一个 $S^1$-主丛;
                \item $\mZ_2$ 通过对径映射作用在 $S^n$ 上,
                商空间是 $\mR P^n$. 于是 $\pi:S^n\rightarrow\mR P^n$ 
                给出一个 $\mZ_2$-主丛;
                \item 设 $H$ 为李群 $G$ 的闭子群, $H$ 通过右乘作用在 $G$ 上, 
                则李群中的一个重要定理告诉我们 $G/H$ 有光滑流形结构,
                且商映射是光滑的. 于是 $\pi:G\rightarrow G/H$ 给出一个 
                $H$-主丛. 
                \item 设 $E$ 为 $M$ 上的秩 $k$ 向量丛, 则 $E$ 的所有标架构成
                $M$ 上的 $GL(k,\mR)$-主丛, 称为 $E$ 的标架丛, 
                记为 $\mathrm{Fr}(E)$. 
            \end{itemize}
        \end{example}
        \begin{remark}
            从第一个定义可以看出通过李群作用能构造大量的主丛, 
            事实上, 若 $G$ 在流形 $M$ 上的作用是自由且 {\rm proper} 的, 
            则微分流形的定理告诉我们商空间 $M/G$ 是流形, 且商映射是光滑的,
            于是 $\pi:M\rightarrow M/G$ 给出一个 $G$-主丛. 
        \end{remark}